\section{Local Asymmetry in the Dalitz Plane}
\label{sec:dalitz}
\label{sec:bayes_dalitz}
\emph{In short: the asymmetry is not uniform over the Dalitz plane ($\chi^2=80.8$ for $31$ degrees of freedom); it is concentrated at low $m^2_{\mathrm{lo}}$ outside the $\phi(1020)$ and, in the region of the published analysis of the same data, it reproduces the published value.}
To test whether the asymmetry is uniform over phase space or concentrated in specific regions, each candidate was placed in the Dalitz plane using the two opposite-sign $K^+K^-$ invariant masses $m^2(H1,H2)$ and $m^2(H1,H3)$ (Sec.~\ref{sec:data_presel}), sorted into $m^2_{\mathrm{lo}}$ and $m^2_{\mathrm{hi}}$ to remove the $H2/H3$ labeling ambiguity. The plane $[0,16]\times[0,30]\,\GeV^2$ in $(m^2_{\mathrm{lo}},m^2_{\mathrm{hi}})$ was divided into an $8\times8$ grid, fixed before looking at the asymmetries. In each bin with at least 20 candidates (32 of 64), the local asymmetry $A_i=(N_i^--N_i^+)/N_i$, with binomial uncertainty $\sigma(A_i)=\sqrt{(1-A_i^2)/N_i}$, was compared with the global counting asymmetry $A_{\mathrm{glob}}=-0.043$ through $S_i=(A_i-A_{\mathrm{glob}})/\sigma(A_i)$; a comparison with zero would instead mostly reflect the global asymmetry itself.

\begin{figure*}[t]
\centering
\includegraphics[width=0.95\textwidth]{figures/bayes_nominal_dalitz_maps.pdf}
\caption{Fixed $8\times8$ Dalitz-plot binning, all candidates in the fit window: (a) number of candidates per bin, (b) local asymmetry $A_i$ for bins with $\geq20$ candidates, and (c) local significance with respect to the global asymmetry, $(A_i-A_{\mathrm{glob}})/\sigma(A_i)$. The two bins with $|S_i|>3$ are outlined. Dashed lines: the squared masses of the $D^0$, $J/\psi$ and $\chi_{c0}$, in either opposite-sign pair.}
\label{fig:bayes_dalitz}
\end{figure*}

Two of the 32 bins have $|S_i|>3$, adjacent to each other at $m^2_{\mathrm{lo}}<4\,\GeV^2$ and $7.5<m^2_{\mathrm{hi}}<11.25\,\GeV^2$ ($S_i=-3.7$ and $-3.0$), with $A_i$ more negative than the global value; the bin below the first one ($m^2_{\mathrm{lo}}<2$, $3.75<m^2_{\mathrm{hi}}<7.5\,\GeV^2$) has $S_i=-2.98$, just under the threshold (Fig.~\ref{fig:bayes_dalitz}); about $0.09$ would be expected if the asymmetry were uniform, with a probability of $3\times10^{-3}$ for two or more. Because the bins are examined together this is a look-elsewhere situation: the number refers to any two of the $32$ bins, not to the two found, and the global homogeneity test that follows does not depend on which bins are singled out. The largest positive deviations, close to $+3$, lie at high $m^2_{\mathrm{hi}}$. A global test of homogeneity, $\chi^2=\sum_i S_i^2=80.8$ for 31 degrees of freedom ($p=3\times10^{-6}$), confirms that the local asymmetry is not uniform. Restricting the sample to the signal window ($|m-\mu|<2\sigma$, $n=15{,}130$, where the fitted model gives a signal purity of $\approx$96\%) gives the same picture -- three bins with $|S_i|>3$, $\chi^2=95.6/28$ -- so the structure is not produced by the combinatorial background. Nor is it produced by the charm contributions: with the $D^0$ veto two bins still exceed $3\sigma$ ($\chi^2=79.2/31$), and with all three vetoes one remains ($\chi^2=74.4/31$, $p=2\times10^{-5}$). It does not depend on the binning either: $6\times6$, $10\times10$ and $12\times12$ grids give $\chi^2=81.0/19$, $97.1/43$ and $120.0/60$ ($p\leq7\times10^{-6}$).

Two null tests check the procedure itself. Shuffling the charges among the candidates removes any local asymmetry while keeping the occupancy of the bins and the global asymmetry of the data: over $2000$ permutations the $8\times8$ grid gives $\chi^2=31.4\pm8.1$, as expected for 31 degrees of freedom, and no permutation reaches the observed value (permutation $p<5\times10^{-4}$; $p=1\times10^{-3}$ for the $12\times12$ grid). The phase-space simulation of the record, without asymmetry by construction, gives $\chi^2=40.0/30$ ($p=0.11$) with the same procedure on the whole sample; on $500$ subsamples of the size of the data the mean $\chi^2/\mathrm{ndof}$ is $1.16$ and $13\%$ have $p<0.05$, so the $\chi^2$ $p$-values are slightly optimistic (the binomial uncertainty is evaluated at the observed $A_i$). The permutation test, whose null distribution is exact, gives the same conclusion: the non-uniformity is a property of the data.

The structure does not depend on the magnet polarity either. The left-right part of the detection asymmetry changes sign when the polarity is reversed, so a local asymmetry produced by it would appear with opposite deviations in the two samples. In the region of the two significant bins ($m^2_{\mathrm{lo}}<4\,\GeV^2$, $7.5<m^2_{\mathrm{hi}}<11.25\,\GeV^2$) the counting asymmetry is $-0.129\pm0.023$ with the magnet down and $-0.126\pm0.028$ with the magnet up, although the global asymmetries of the two samples differ ($-0.051$ and $-0.032$), and the two maps agree bin by bin ($\chi^2=23.0$ for 29 bins, $p=0.78$). This split does not test the part of the detection asymmetry that does not change sign with the polarity, the different interaction of $K^+$ and $K^-$ with the detector material.

\begin{figure*}[t]
\centering
\includegraphics[width=0.95\textwidth]{figures/ml_nominal_dalitz_region.pdf}
\caption{(a) Ordered Dalitz plot of the candidates in the fit window, with the squared masses of the $\phi(1020)$, $D^0$, $J/\psi$ and $\chi_{c0}$ (dashed) and the region of the published analysis of the same data~\cite{lhcb_cpv_2013} (box). (b) Mass spectra of $B^+$ and $B^-$ candidates in that region, same binning: the $B^-$ peak is visibly smaller.}
\label{fig:dalitz_region}
\end{figure*}

The ordered Dalitz plot (Fig.~\ref{fig:dalitz_region}a) shows the two-body resonances of the decay: the narrow $\phi(1020)$ band at the kinematic edge, $m^2_{\mathrm{lo}}\approx1.04\,\GeV^2$, and the charm bands of the $D^0$, $J/\psi$ and $\chi_{c0}$. The negative local asymmetry is not produced by the $\phi(1020)$: its band ($|m_{\mathrm{lo}}-m_\phi|<10\,\MeV$, $1964$ candidates) has $A_{\mathrm{count}}=-0.004\pm0.023$, while the significant bins without it have $-0.145\pm0.019$. The asymmetry sits in the low-mass $K^+K^-$ region just above the $\phi$, which is where the published analysis of the same data located it, defining the region $1.2<m^2_{\mathrm{lo}}<2.0\,\GeV^2$, $m^2_{\mathrm{hi}}<15\,\GeV^2$ and measuring there $\Araw=-0.239\pm0.020$~\cite{lhcb_cpv_2013}. Since that region is fixed by the paper and not chosen on our data, it allows an unbiased comparison: fitting the $2146$ candidates in it with the nominal model gives $\Araw=-0.234\pm0.025$ (Bayesian; the chain mixes slowly in the poorly constrained second background exponential, $N/\tau=15$) and $-0.237\pm0.025$ (frequentist, $9.4\sigma$ from zero), reproducing the published value; the mass spectra of the two charges in the region (Fig.~\ref{fig:dalitz_region}b) show the difference directly. Three caveats apply: the binning is coarse and was not optimized; the asymmetries are computed on all candidates, without the sWeight-based background subtraction that a full analysis would use (Sec.~\ref{sec:limitations}); and, since $\Araw$ is not corrected for detection asymmetries, local variations can still be partly instrumental -- the $K^+$/$K^-$ interaction asymmetry depends on the kaon momenta, which are correlated with the position in the Dalitz plane, and is not removed by the polarity split. Appendix~\ref{sec:ml_kmeans} revisits the same question with an adaptive, data-driven partition ($K$-means).




